% =============================================================================
%  O tempo de reacao do autoescalamento e uma propriedade da plataforma?
%  PSI5120 - Topicos em Computacao em Nuvem (2026) - Trabalho Final
% -----------------------------------------------------------------------------
%  FORMATO: IEEE conference (IEEEtran), duas colunas, EM PORTUGUES.
%
%  IDIOMA: esta e a versao em portugues. A versao em ingles esta em
%  artigo-final-en.tex. As duas leem as MESMAS tabelas geradas pela analise, em
%  arquivos que diferem apenas no idioma dos rotulos (dados/tabelas-pt.tex e
%  dados/tabelas-en.tex), de modo que nao podem divergir numericamente.
%
%  COMO COMPILAR (a partir da pasta artigo/):
%      bash ../scripts/12-prep-artigo.sh     % copia os CSVs da analise
%      pdflatex artigo-final.tex
%      pdflatex artigo-final.tex             % 2a passada: referencias cruzadas
%
%  OU, automaticamente:  bash scripts/17-compilar-artigo.sh
%
%  FIGURAS: nao ha imagens intermediarias. Todos os graficos sao gerados pelo
%  pgfplots lendo DIRETAMENTE os CSVs produzidos por scripts/11-analyze.py, e
%  as tabelas de dados vem de dados/tabelas-pt.tex, gerado pelo mesmo script.
%  Consequencia pratica: e impossivel um numero do artigo discordar do dado
%  medido -- se a campanha for repetida, o PDF se atualiza sozinho.
% =============================================================================

\documentclass[conference]{IEEEtran}

\usepackage[utf8]{inputenc}
\usepackage[T1]{fontenc}
\usepackage[brazil]{babel}

\usepackage{graphicx}
\usepackage{booktabs}
\usepackage{array}
\usepackage{amsmath}
\usepackage{url}
\usepackage{cite}
\usepackage[hidelinks]{hyperref}

\usepackage{pgfplots}
\usepackage{pgfplotstable}
\pgfplotsset{compat=1.18}
\usepgfplotslibrary{fillbetween}
\usetikzlibrary{positioning, arrows.meta, fit, backgrounds}

% O IEEEtran desativa a hifenizacao em alguns contextos; o portugues tem
% palavras longas e sofre mais em colunas estreitas, entao reforcamos os pontos
% de divisao das palavras que mais aparecem no texto.
\hyphenation{Kuber-netes auto-esca-la-men-to esca-lo-na-men-to re-pro-du-ti-bi-li-da-de
             in-fra-es-tru-tu-ra dis-po-ni-bi-li-da-de ins-tru-men-ta-cao
             sin-cro-ni-za-cao es-ta-bi-li-za-cao per-mu-ta-cao}

% Paleta unica para os dois ambientes, usada em todas as figuras.
\definecolor{corLocal}{RGB}{196, 94, 32}
\definecolor{corNuvem}{RGB}{31, 94, 158}
\definecolor{corNeutra}{RGB}{110, 110, 110}

% Estilo comum dos graficos: sem molduras superfluas, grade discreta.
\pgfplotsset{
  estiloBase/.style={
    width=\linewidth, height=5.2cm,
    tick align=outside, tick pos=left,
    grid=major, grid style={line width=0.3pt, draw=black!15},
    every axis plot/.append style={line width=1pt},
    legend style={font=\scriptsize, draw=black!25, fill=white, fill opacity=0.9,
                  text opacity=1, cells={anchor=west}},
    label style={font=\small}, tick label style={font=\scriptsize},
    title style={font=\small\bfseries},
  }
}

\begin{document}

% =============================================================================
\title{O tempo de reação do autoescalamento é uma\\
propriedade da plataforma? Um estudo de medição\\
replicado do HPA do Kubernetes em Minikube e Amazon EKS}

\author{
\IEEEauthorblockN{José Vitor Feitosa de Andrade}
\IEEEauthorblockA{Escola Politécnica\\
Universidade de São Paulo\\
São Paulo, Brasil\\
NUSP 13682041\\
\texttt{jose.vitor.0905@gmail.com}}
\and
\IEEEauthorblockN{Fernando Frederice Miqueletti}
\IEEEauthorblockA{Escola Politécnica\\
Universidade de São Paulo\\
São Paulo, Brasil\\
NUSP 12544340\\
\texttt{fernando.miqueletti@usp.br}}
}

\maketitle

% =============================================================================
\begin{abstract}
Comparações de infraestrutura de nuvem feitas por profissionais costumam se
apoiar em uma única execução por configuração. Este trabalho mostra, num cenário
em que replicar é barato, o quanto esse desenho pode falhar --- e mostra isso
corrigindo uma conclusão nossa. Um estudo anterior implantou uma carga de
trabalho Kubernetes idêntica, com autoescalamento horizontal de Pods, em um
\emph{cluster} Minikube local e no serviço gerenciado Amazon EKS, executou cada
implantação uma vez por integrante, observou que o ambiente mais rápido mudava
entre as execuções e concluiu que o tempo de reação não era atribuível à
plataforma. Este artigo mantém inalterados a implantação, a carga e o protocolo,
e altera apenas o desenho experimental: vinte repetições não supervisionadas por
ambiente, com a métrica primária derivada exclusivamente de instantes atribuídos
pelo \emph{cluster} --- o \texttt{lastScaleTime} do próprio controlador e o
\texttt{startedAt} do contêiner gerador de carga --- de modo que nem o relógio do
observador nem sua cadência de consulta entram na medição. Nas quarenta
repetições, sem nenhuma perda, o tempo de reação foi de 40,8\,s no ambiente local
contra 24,7\,s na nuvem (Mann--Whitney $U=52,0$, $p=1\times10^{-4}$;
$\delta$ de Cliff $=0{,}74$): existe um efeito de plataforma real, que o desenho
anterior não tinha como resolver. E as mesmas distribuições explicam por que ele
não tinha. Uma repetição sorteada de cada ambiente reproduz a ordenação correta
em apenas 87\,\% dos casos, e as quatro observações relatadas anteriormente caem
nos percentis 82, 10, 45 e 100 das distribuições aqui medidas --- sorteios comuns
que por acaso se inverteram. Registrar covariáveis junto de cada repetição
permite testar explicações candidatas em vez de supô-las: nenhuma das cinco
--- latência de escrita no \emph{control plane}, sua parcela de armazenamento,
partida do gerador, utilização de pico ou deriva ao longo da campanha --- prediz
quais repetições locais foram lentas, apesar de a latência de escrita local
variar de 114\,ms a 2534\,ms. Já tudo o que é fixado por declaração se reproduziu
com variação de poucos por cento. Concluímos que o resultado lógico do
autoescalamento pode ser validado num \emph{cluster} local, que sua latência não
pode, e que a aritmética de uma comparação de execução única com $A_{12}=0{,}87$
deveria desencorajar a prática muito além deste experimento.

\vspace{0.6em}
\noindent\textit{Abstract}---Practitioner comparisons of cloud infrastructure
routinely rest on a single execution per configuration. We show, in a setting
where replication is cheap, how badly that design can fail --- and we show it by
correcting a conclusion of our own. An earlier study deployed an identical
Kubernetes workload with the Horizontal Pod Autoscaler on a local Minikube
cluster and on managed Amazon EKS, executed each deployment once per author, and
observed that the faster environment changed between executions; it concluded
that autoscaling reaction time was not attributable to the platform. This paper
keeps the deployment, the workload and the load protocol unchanged, and changes
only the experimental design: twenty unattended repetitions per environment, with
the primary metric derived exclusively from timestamps assigned by the cluster.
Across forty repetitions with no attrition, reaction time averaged 40.8\,s
locally against 24.7\,s in the cloud (Mann--Whitney $U=52.0$,
$p=1\times10^{-4}$; Cliff's $\delta=0.74$): a real platform effect that the
earlier design could not resolve. A single repetition drawn from each environment
reproduces the correct ordering only 87\,\% of the time, and the four previously
reported observations sit at the 82nd, 10th, 45th and 100th percentiles of the
distributions measured here. None of five recorded covariates predicts which
local repetitions were slow. Everything fixed by declaration replicated to within
a few percent.
\end{abstract}

\begin{IEEEkeywords}
Computação em nuvem, Kubernetes, autoescalamento, \emph{Horizontal Pod
Autoscaler}, elasticidade, reprodutibilidade, metodologia de medição, Amazon EKS.
\end{IEEEkeywords}

% =============================================================================
\section{Introdução}
\label{sec:intro}

A elasticidade rápida é uma das cinco características essenciais da computação em
nuvem na definição do NIST~\cite{nist}: recursos devem ser provisionados e
liberados automaticamente, de forma proporcional à demanda. No Kubernetes, o
mecanismo que entrega essa propriedade na camada de aplicação é o
\emph{Horizontal Pod Autoscaler} (HPA)~\cite{k8shpa}, uma malha de controle que
periodicamente compara uma métrica observada com uma meta e ajusta o número de
réplicas de uma carga de trabalho.

Quem escolhe entre um \emph{cluster} local de desenvolvimento e um serviço
gerenciado de nuvem pergunta, naturalmente, se o autoescalamento \emph{se
comporta} igual nos dois. A pergunta tem duas metades distintas, e confundi-las é
fonte de uma quantidade surpreendente de mal-entendidos. A primeira metade é o
resultado \emph{lógico}: para quantas réplicas o controlador converge, e se a
utilização observada se estabiliza na meta configurada. A segunda é o resultado
\emph{temporal}: quanto tempo o sistema leva para reagir a um pico de demanda.

Nosso próprio estudo anterior sobre essa questão ilustra a armadilha. No trabalho
intermediário desta disciplina, implantamos uma carga de trabalho idêntica --- a
imagem oficial \texttt{hpa-example}, com manifestos, \emph{requests},
\emph{limits}, meta de utilização e política de escalonamento idênticos --- em um
\emph{cluster} Minikube local e em um \emph{cluster} gerenciado Amazon EKS, e
cada integrante do grupo executou o experimento completo de forma independente. O
resultado lógico replicou perfeitamente: as quatro implantações convergiram para
seis ou sete réplicas, com a utilização média de CPU estabilizando na meta de
50\,\%. O resultado temporal não replicou. Na primeira execução, o EKS reagiu
cerca de quatro vezes mais rápido que o Minikube (15\,s contra 62\,s); na
segunda, realizada em outra máquina e em outra conta AWS com o mesmo protocolo, a
ordenação \emph{se inverteu}, com o Minikube reagindo cerca de duas vezes mais
rápido que o EKS.

Com uma execução por ambiente, a análise tinha de parar ali. Duas observações
bastam para mostrar que um resultado não é estável, mas não bastam para dizer
qualquer coisa sobre a distribuição subjacente: não há média, não há dispersão,
não há intervalo de confiança e não há teste. A conclusão natural --- ``a
diferença não é atribuível à plataforma'' --- era uma afirmação sobre a
\emph{ausência} de consistência, não uma caracterização do fenômeno.

Este artigo fecha essa lacuna. Mantemos deliberadamente inalterados a
implantação, a carga e o protocolo do estudo anterior, e mudamos apenas o desenho
experimental: em vez de uma execução por ambiente, conduzimos uma campanha
automatizada e não supervisionada de execuções repetidas, e tratamos o tempo de
reação como uma variável aleatória em vez de um número. Revisamos também a
\emph{procedência} de cada instante medido, derivando a métrica primária
exclusivamente de carimbos de tempo atribuídos pelo próprio \emph{cluster}, de
modo que nem o relógio local nem a cadência de consulta do observador possam
contaminá-la.

As questões de pesquisa são:

\begin{itemize}
\item[\textbf{QP1}] Qual é a \emph{distribuição} do tempo de reação do HPA em
cada ambiente, e ela difere entre um \emph{cluster} local e um gerenciado por
mais do que ruído de amostragem?
\item[\textbf{QP2}] Os dois ambientes diferem em \emph{dispersão} --- isto é, o
ambiente local não é apenas mais lento ou mais rápido, mas menos
\emph{previsível}?
\item[\textbf{QP3}] Dadas as distribuições medidas, qual é a probabilidade de um
desenho de execução única --- o desenho do estudo anterior, e da maior parte dos
relatos de profissionais --- conduzir a cada conclusão possível?
\item[\textbf{QP4}] Quais partes do comportamento de autoescalamento \emph{são}
reprodutíveis entre ambientes e, portanto, podem legitimamente ser validadas num
\emph{cluster} local?
\end{itemize}

A contribuição é tripla: um método de instrumentação que torna a latência de
autoescalamento mensurável sem confiar no relógio do observador
(Seção~\ref{sec:instrumentacao}); uma caracterização empírica da distribuição do
tempo de reação nos dois ambientes, com tratamento não paramétrico completo
(Seção~\ref{sec:resultados}); e uma resposta quantificada a quão enganoso é um
experimento de execução única neste cenário
(Seção~\ref{sec:execucao-unica}) --- uma questão que generaliza bem além do HPA.

% =============================================================================
\section{Fundamentação}
\label{sec:background}

\subsection{A malha de controle do HPA}

O HPA é uma malha de controle fechada executada pelo
\texttt{kube-controller-manager} em um período fixo
(\texttt{-{}-horizontal-pod-autoscaler-sync-period}, 15\,s por padrão). A cada
iteração ele lê o valor corrente da métrica configurada na \emph{Metrics API} e
calcula o número desejado de réplicas como~\cite{k8salgo}

\begin{equation}
\label{eq:hpa}
r_{\text{desejado}} = \left\lceil\, r_{\text{atual}} \times
\frac{m_{\text{atual}}}{m_{\text{meta}}} \,\right\rceil
\end{equation}

\noindent onde $m$ é, em nossa configuração, a utilização média de CPU dos Pods
expressa como fração de seu \emph{request} de CPU. O resultado é limitado ao
intervalo $[\texttt{minReplicas}, \texttt{maxReplicas}]$ e ainda restringido pela
cláusula \texttt{behavior}, que define janelas de estabilização e políticas de
taxa de forma independente para subida e descida.

A Equação~\eqref{eq:hpa} é puramente declarativa: não contém nenhum termo que
dependa da infraestrutura. É precisamente por isso que se espera que o resultado
\emph{lógico} do autoescalamento replique entre ambientes, e isso delimita a
questão que este artigo aborda --- tudo o que pode diferir entre um
\emph{cluster} local e um gerenciado está na \emph{temporização} da malha, não em
sua aritmética.

\subsection{Onde a latência pode entrar}
\label{sec:latencia}

Entre o instante em que uma carga se torna limitada por CPU e o instante em que
uma nova réplica existe, acumulam-se vários atrasos independentes:

\begin{enumerate}
\item \textbf{Aquisição da métrica.} O cAdvisor do kubelet amostra a CPU dos
contêineres periodicamente; o \emph{Metrics Server}~\cite{metricsserver} coleta
dos kubelets em seu próprio período (15\,s por padrão) e serve uma janela móvel
curta. Um pico de CPU, portanto, só se torna visível ao HPA depois de até um
período completo de coleta, mais a largura da janela sobre a qual a taxa é
calculada.
\item \textbf{Sincronização do controlador.} O HPA avalia a
Equação~\eqref{eq:hpa} uma vez por período de sincronização. No pior caso, a
métrica se torna visível imediatamente após uma avaliação e espera mais um
período inteiro.
\item \textbf{Estabilização e políticas de taxa.} Com
\texttt{scaleUp.stabilizationWindowSeconds} igual a zero, como em nossa
configuração, nenhum atraso adicional é imposto na subida; a política percentual
ainda limita o quanto a contagem de réplicas pode crescer por intervalo.
\item \textbf{Atuação.} Uma vez escalado o Deployment, o escalonador precisa
alocar os novos Pods e o kubelet precisa iniciá-los --- um atraso que depende do
cache de imagens, da capacidade do nó e, em um \emph{cluster} gerenciado, do
escalonamento entre nós.
\end{enumerate}

Apenas os itens 1 e 2 precedem a \emph{decisão}. Essa distinção importa para a
definição da métrica na Seção~\ref{sec:instrumentacao}: medimos o tempo até a
\emph{decisão} do controlador, não até os novos Pods estarem prontos, e
reportamos os atrasos de atuação separadamente. A quantização implicada pelos
itens 1 e 2 também explica por que não se deve esperar que o tempo de reação seja
normalmente distribuído --- é uma quantidade limitada inferiormente e
discretizada, o que motiva o tratamento não paramétrico da
Seção~\ref{sec:estatistica}.

\subsection{Reprodutibilidade em medição de sistemas}

A preocupação metodológica deste artigo não é nova. É bem estabelecido que
medições de desempenho em infraestrutura compartilhada estão sujeitas a
variabilidade grande o bastante para inverter conclusões, e que experimentos de
execução única são correspondentemente frágeis~\cite{mytkowicz, hoefler}. A
recomendação prática --- relatar distribuições e tamanhos de efeito em vez de
comparações pontuais, e usar testes não paramétricos quando a distribuição
subjacente é desconhecida --- é padrão na literatura de metodologia
experimental~\cite{arcuri}. O que este artigo acrescenta é uma instância concreta
no domínio de orquestração em nuvem, onde o custo de replicar é quase nulo, a
variabilidade é alta, e a literatura --- inclusive nosso próprio trabalho
anterior --- rotineiramente relata execuções únicas.

% =============================================================================
\section{O trabalho anterior e sua limitação}
\label{sec:anterior}

O trabalho intermediário implantou a carga descrita na
Seção~\ref{sec:desenho} em Minikube e em Amazon EKS. Cada um dos dois autores
executou o experimento completo de forma independente, em máquinas e contas AWS
distintas, aplicando manifestos idênticos byte a byte e o mesmo protocolo de
carga. Os resultados estão resumidos na Tabela~\ref{tab:anterior}.

\begin{table}[!t]
\centering
\caption{Resultados do estudo anterior, de execução única. O resultado lógico
replicou; o temporal não.}
\label{tab:anterior}
% \scriptsize e rotulos curtos: cinco colunas nao cabem em uma coluna do IEEE
% com \footnotesize (overfull de 12,6 pt).
\scriptsize
\begin{tabular}{@{}lrrrr@{}}
\toprule
& \multicolumn{2}{c}{\textbf{Execução E1}} & \multicolumn{2}{c}{\textbf{Execução E2}}\\
\cmidrule(lr){2-3}\cmidrule(lr){4-5}
\textbf{Métrica} & Minikube & EKS & Minikube & EKS\\
\midrule
Réplicas no regime & 6 & 6 & 6 & 7\\
CPU no regime & 51\,\% & 47\,\% & $\sim$50\,\% & $\sim$50\,\%\\
CPU de pico & 153\,\% & 237\,\% & 166\,\% & n/d\\
\textbf{Tempo de reação} & \textbf{62\,s} & \textbf{15\,s} & \textbf{$\sim$33\,s} & \textbf{$\sim$60\,s}\\
Descida ao mínimo & 7,6\,min & 6,8\,min & 13--14\,min & 9--10\,min\\
\bottomrule
\end{tabular}
\end{table}

Dois aspectos da Tabela~\ref{tab:anterior} motivam este artigo. Primeiro, as
linhas que descrevem o resultado \emph{lógico} são estáveis nas quatro
implantações. Segundo, a linha do tempo de reação não é apenas ruidosa --- a
\emph{ordenação} dos dois ambientes se inverte entre E1 e E2. Se apenas E1
tivesse sido realizada, o estudo teria concluído, com dados internamente
consistentes e uma explicação mecanística plausível, que o \emph{control plane}
gerenciado reage quatro vezes mais rápido. Se apenas E2 tivesse sido realizada,
teria concluído o oposto.

O desenho do estudo anterior deixou também três fragilidades de instrumentação
que este artigo remove. O tempo de reação era calculado a partir do relógio de
parede no momento em que o operador lançava o gerador de carga, o que inclui o
escalonamento e a partida do Pod --- uma quantidade que difere
sistematicamente entre um \emph{cluster} local de nó único com cache de imagens
quente e um \emph{cluster} gerenciado de dois nós. O período de amostragem era de
5\,s, grosseiro em relação ao período de 15\,s do controlador. Por fim, as duas
execuções usaram instrumentação heterogênea: E1 dispunha de um amostrador
automatizado, enquanto E2 se baseou em observação direta, de modo que parte da
diferença observada poderia ser de medição, não de fenômeno.

% =============================================================================
\section{Desenho experimental}
\label{sec:desenho}

\subsection{O que é mantido constante}

O sistema implantado é idêntico byte a byte ao do estudo anterior, e idêntico
entre os dois ambientes. A Tabela~\ref{tab:config} lista a configuração.
Mantê-la fixa é deliberado: torna a campanha diretamente comparável tanto às
execuções únicas anteriores quanto entre ambientes, de modo que qualquer
diferença observada seja atribuível ao ambiente e não à especificação.

\begin{table}[!t]
\centering
\caption{Configuração mantida idêntica entre ambientes e campanhas.}
\label{tab:config}
\footnotesize
% Coluna de valor como p{}: a linha do gerador de carga nao cabe em uma celula
% de largura natural (overfull de 40 pt com especificacao "ll").
\begin{tabular}{@{}p{0.52\linewidth}p{0.42\linewidth}@{}}
\toprule
\textbf{Parâmetro} & \textbf{Valor}\\
\midrule
Imagem da aplicação & \texttt{registry.k8s.io/\allowbreak hpa-example}\\
\emph{Request} / \emph{limit} de CPU & 200\,m / 500\,m\\
\emph{Request} / \emph{limit} de memória & 64\,Mi / 128\,Mi\\
Versão da API do HPA & \texttt{autoscaling/v2}\\
Métrica & utilização de CPU, média\\
Meta & 50\,\% do \emph{request}\\
\texttt{minReplicas} / \texttt{maxReplicas} & 1 / 10\\
\texttt{scaleUp.\allowbreak stabilizationWindow} & 0\,s\\
Política de \texttt{scaleUp} & 100\,\% a cada 15\,s\\
\texttt{scaleDown.\allowbreak stabilizationWindow} & 300\,s\\
Política de \texttt{scaleDown} & 50\,\% a cada 60\,s\\
Gerador de carga & um Pod \texttt{busybox:1.28}, laço sequencial\\
Duração da carga por repetição & 240\,s\\
\bottomrule
\end{tabular}
\end{table}

O gerador de carga merece comentário. É um único Pod executando
\texttt{while sleep 0.01; do wget -q -O- http://php-apache; done} --- um laço de
requisições sequencial e de uma única linha de execução. É um gerador de carga
fraco, e o estudo anterior o identificou como o fator que impediu qualquer dos
\emph{clusters} de alcançar o teto de dez réplicas. Nós o mantemos inalterado,
porque o objeto de estudo aqui é a \emph{reação} do controlador ao início da
carga, não o limite de capacidade dos \emph{clusters}; substituir o gerador
tornaria a campanha incomparável a E1 e E2, que é justamente a comparação sobre a
qual o artigo se constrói. A Seção~\ref{sec:ameacas} retoma as consequências
dessa escolha.

\subsection{Ambientes}

O ambiente local é um \emph{cluster} Minikube de nó único sobre o \emph{driver}
Docker, executando em um hospedeiro Windows~11 com WSL2. O ambiente de nuvem é um
\emph{cluster} gerenciado Amazon EKS criado com \texttt{eksctl}~\cite{eksctl},
com um \emph{node group} gerenciado de duas instâncias \texttt{t3.small} em
\texttt{us-east-1} --- a mesma topologia do estudo anterior. As características
efetivamente reportadas por cada \emph{cluster} em tempo de execução, e não as
solicitadas na criação, estão na Tabela~\ref{tab:ambientes}; a distinção importa,
como a Seção~\ref{sec:ameacas} explica.

% Valores lidos dos proprios clusters em tempo de execucao; fonte:
% dados/<ambiente>/ambiente.txt.
\begin{table}[!t]
\centering
\caption{Características dos ambientes conforme reportadas por cada
\emph{cluster} em tempo de execução, e não conforme solicitadas na criação.}
\label{tab:ambientes}
\footnotesize
\begin{tabular}{@{}p{0.30\linewidth}p{0.30\linewidth}p{0.30\linewidth}@{}}
\toprule
\textbf{Característica} & \textbf{Minikube} & \textbf{Amazon EKS}\\
\midrule
Kubernetes (\emph{API server}) & v1.35.1 & v1.34.11-eks\\
Nós & 1 (\emph{control plane} e trabalhador) & 2 trabalhadores\\
CPU alocável por nó & 4000\,m & 1930\,m\\
Memória alocável por nó & 7,8\,GiB & 1,4\,GiB\\
Tipo de instância & n/d (contêiner em Windows\,11 + WSL2) & \texttt{t3.small}\\
SO do nó & Debian 12 & Amazon Linux 2023\\
\emph{Runtime} de contêiner & Docker 29.2.1 & containerd 2.2.7\\
Zonas de disponibilidade & n/d & \texttt{us-east-1b}, \texttt{us-east-1f}\\
\emph{Metrics Server} & \emph{addon} do Minikube & \emph{addon} gerenciado do EKS\\
Desvio de relógio cliente--\emph{cluster} & 0\,s & $-2$\,s\\
\bottomrule
\end{tabular}
\end{table}

Duas entradas da Tabela~\ref{tab:ambientes} merecem comentário. O \emph{cluster}
local foi criado com \texttt{-{}-cpus=2 -{}-memory=4096}, mas reporta 4000\,m de
CPU alocável e 7,8\,GiB de memória: o \emph{driver} Docker Desktop não aplica
esses limites, comportamento já observado no estudo anterior e aqui confirmado. O
nó local é, portanto, \emph{maior} que os dois nós de nuvem somados (4000\,m
contra $2\times1930$\,m), o que importa na interpretação dos resultados ligados a
capacidade. O desvio de relógio é reportado apenas por completude; a métrica
primária é uma diferença entre dois carimbos de tempo atribuídos pelo
\emph{cluster} e não depende do relógio do cliente.

\subsection{Protocolo de cada repetição}

Uma campanha consiste em $N$ repetições do protocolo abaixo, executadas sem
supervisão por \texttt{scripts/10-run-campaign.sh}:

\begin{enumerate}
\item \textbf{Reinício.} Qualquer gerador de carga é removido e o Deployment é
escalado para uma réplica.
\item \textbf{Verificação do estado de base.} A repetição não começa até que o
HPA reporte uma réplica desejada, uma réplica pronta e utilização de CPU
observada igual ou inferior a 10\,\% por três amostras consecutivas. Como
\texttt{scaleUp.stabilizationWindowSeconds} é zero, o controlador desfaz uma
redução manual enquanto o \emph{Metrics Server} ainda reporta a CPU decaindo da
repetição anterior; o laço, por isso, reaplica a escala até que a CPU medida
tenha de fato decaído. Sem essa etapa, as repetições não seriam independentes.
\item \textbf{Amostragem.} Um amostrador registra o estado do HPA e do Deployment
a cada 2\,s, com compensação de deriva para que o intervalo efetivo não cresça
quando chamadas individuais à API demoram mais.
\item \textbf{Carga.} O Pod gerador é submetido. O $T_0$ da repetição é o
instante em que seu contêiner \emph{efetivamente começou} a executar, não o
instante da submissão (Seção~\ref{sec:instrumentacao}).
\item \textbf{Detecção.} A primeira decisão de escala do controlador é detectada
consultando \texttt{hpa.status.lastScaleTime} uma vez por segundo e tomando o
primeiro valor que seja, ao mesmo tempo, diferente do valor anterior à carga e
não anterior a $T_0$.
\item \textbf{Manutenção e encerramento.} A carga é mantida por 240\,s e o
gerador é removido.
\item \textbf{Descida.} Nas três primeiras repetições de cada campanha a descida
natural é observada até o fim; nas demais, o retorno a uma réplica é forçado,
para manter a campanha em duração limitada. A janela de estabilização de 300\,s
faz uma descida natural completa custar cerca de dez minutos por repetição, o que
mais que dobraria a campanha e, no ambiente de nuvem, seu custo.
\end{enumerate}

\subsection{Instrumentação: a procedência de cada instante}
\label{sec:instrumentacao}

A escolha metodológica central deste artigo é que a métrica primária é calculada
a partir de dois carimbos de tempo \emph{atribuídos pelo cluster}, nunca pelo
cliente que mede:

\begin{itemize}
% O caminho completo do campo nao cabe em uma coluna de ~8,8 cm e o \texttt
% nao oferece pontos de quebra; \allowbreak depois de cada ponto resolve sem
% alterar o texto impresso.
\item $T_0$ é \texttt{pod.\allowbreak status.\allowbreak
containerStatuses[0].\allowbreak state.\allowbreak running.\allowbreak
startedAt} do gerador de carga --- carimbado pelo kubelet quando o contêiner
começou a executar.
\item $T_1$ é \texttt{hpa.status.lastScaleTime} --- carimbado pelo próprio
controlador do HPA quando ele agiu.
\item \textbf{Tempo de reação} $:= T_1 - T_0$.
\end{itemize}

Como ambos os instantes se originam dentro do \emph{cluster}, a medição é imune
ao relógio local, ao tempo de ida e volta na rede e ao instante em que o
amostrador por acaso perguntou. Isso não é um refinamento cosmético. No estudo
anterior, $T_0$ era o relógio de parede do operador na submissão, o que embute
silenciosamente o escalonamento e o \emph{pull} da imagem do Pod no tempo de
reação relatado --- e esse atraso é exatamente o tipo de quantidade que difere
entre um nó local quente e um \emph{cluster} gerenciado de dois nós. Aqui ele é
medido como variável separada (\emph{partida do gerador}, o intervalo da
submissão até o início do contêiner) e reportado de forma independente.

\paragraph{Manter a amostragem simétrica} Um obstáculo prático ameaçou
reintroduzir justamente o defeito que nos propusemos a remover. O kubeconfig
gerado por \texttt{aws eks update-kubeconfig} autentica através de um \emph{exec
credential plugin}: cada invocação do \texttt{kubectl} lança a AWS CLI para
emitir um \emph{token}. No hospedeiro de medição isso custou 2,8\,s por chamada,
de modo que um amostrador configurado para 2\,s entregava, na prática, uma
amostra a cada 5\,s contra o EKS, enquanto entregava 2\,s contra o Minikube. A
métrica primária sobreviveria a isso --- é derivada de carimbos atribuídos pelo
\emph{cluster} e é indiferente a quando o observador pergunta --- mas toda
quantidade derivada da série temporal (utilização de pico, tempo até o regime
permanente) teria sido medida em resoluções diferentes nos dois ambientes. Essa é
exatamente a ameaça de instrumentação heterogênea que o estudo anterior
identificou em si mesmo.

Substituímos então o \emph{exec plugin} por um kubeconfig com um \emph{token}
previamente emitido, renovado em segundo plano a cada sete minutos porque os
\emph{tokens} do EKS têm vida curta. A latência por chamada caiu de 2,8\,s para
0,51\,s, e o período efetivo de amostragem passou a ser constante de 2\,s nos
dois ambientes. Os 0,51\,s residuais são tempo de ida e volta e negociação TLS
até \texttt{us-east-1}, que nenhuma mudança do lado do cliente remove, e são a
razão pela qual os valores de latência do lado do cliente da
Seção~\ref{sec:resultados} devem ser lidos como uma \emph{diferença} entre
escrita e leitura, e não em termos absolutos.

Duas verificações cruzadas acompanham a métrica primária. Primeiro, um tempo de
reação derivado da consulta --- a primeira amostra em que \texttt{desired} excede
um --- é registrado em paralelo; espera-se que seja ligeiramente maior e mais
grosseiro, e qualquer divergência grande indicaria falha de instrumentação.
Segundo, o desvio entre o relógio do cliente e o do \emph{cluster} é medido no
início de cada campanha criando um objeto descartável e comparando o
\texttt{creationTimestamp} atribuído pelo servidor com o relógio local, e é
reportado na Tabela~\ref{tab:ambientes}. Ainda que a métrica primária não dependa
dele, publicar o desvio permite ao leitor limitar o erro de qualquer quantidade
medida do lado do cliente.

% =============================================================================
\section{Métodos estatísticos}
\label{sec:estatistica}

Toda a análise é feita por \texttt{scripts/11-analyze.py}, implementado usando
apenas a biblioteca padrão do Python --- sem numpy, sem scipy. A motivação é
reprodutibilidade: o leitor consegue reexecutar a análise sem instalar nada, e
cada estimador fica visível no código-fonte em vez de escondido atrás de uma
chamada de biblioteca. A semente aleatória é fixa, de modo que os resultados de
\emph{bootstrap} e de permutação são reprodutíveis.

\paragraph{Estatísticas descritivas} Reportamos $n$, média, desvio-padrão
amostral, coeficiente de variação, mínimo, quartis, máximo e desvio absoluto
mediano (MAD) para cada métrica e ambiente.

\paragraph{Estimação por intervalo} Os intervalos de confiança para a média são
calculados por \emph{bootstrap} percentil não paramétrico com $B = 20\,000$
reamostragens, em vez da distribuição $t$. O tempo de reação é limitado
inferiormente por zero e quantizado pelo período de sincronização do controlador
(Seção~\ref{sec:latencia}), de modo que normalidade não é uma suposição
defensável com $n \approx 20$.

\paragraph{Teste de hipótese} Os dois ambientes são comparados com um teste de
Mann--Whitney $U$ bilateral, usando a aproximação normal com correções de
continuidade e de empates. Empates são esperados e precisam ser corrigidos
precisamente porque a medição é quantizada. Como segunda verificação, livre de
suposições, executamos um teste de permutação sobre a diferença de medianas, com
50\,000 reetiquetagens sob a hipótese nula de permutabilidade.

\paragraph{Dispersão} A QP2 trata de variabilidade e não de posição, então a
testamos diretamente: um teste de permutação sobre o logaritmo da razão entre os
dois MADs. Essa é a formulação formal da afirmação de que um ambiente local é
menos \emph{previsível}, o que é distinto de --- e não implicado por --- ser mais
lento em média.

\paragraph{Tamanho de efeito} Reportamos o $\delta$ de Cliff e a estatística
$A_{12}$ de Vargha--Delaney. São os tamanhos de efeito naturais aqui porque
$A_{12}$ tem leitura operacional direta neste experimento: é a probabilidade de
uma repetição sorteada de um ambiente exceder uma repetição sorteada do outro ---
que é exatamente a pergunta que um estudo de execução única responde
implicitamente, e portanto a base da Seção~\ref{sec:execucao-unica}.

% =============================================================================
\section{Resultados}
\label{sec:resultados}

As duas campanhas se completaram sem nenhuma perda: 20 repetições válidas por
ambiente, 2849 e 2988 amostras de série temporal respectivamente, e nenhuma
repetição perdida por \emph{cluster} degradado ou por falha do gerador de carga.
A Tabela~\ref{tab:descritivas} traz as estatísticas descritivas de todas as
métricas; o texto abaixo destaca o que importa para as questões de pesquisa.

% -----------------------------------------------------------------------------
% As tabelas de descritivas e de testes sao geradas por 11-analyze.py.
% Gerado por scripts/11-analyze.py -- nao editar a mao.
% Tabela 1: descritivas por ambiente.
\begin{table*}[!t]
\centering
\caption{Estatísticas descritivas das repetições (IC de 95\,\% da média por bootstrap, $B=20000$).}
\label{tab:descritivas}
\footnotesize
\begin{tabular}{@{}llrrrrrrr@{}}
\toprule
\textbf{Métrica} & \textbf{Ambiente} & $n$ & \textbf{Média} & \textbf{DP} & \textbf{CV} & \textbf{Mediana} & \textbf{IIQ} & \textbf{IC 95\%}\\
\midrule
Tempo de reação do HPA (s) & Minikube & 20 & 40.8 & 18.5 & 0.45 & 33.5 & 30.8--41.5 & [33.5, 49.2]\\
Tempo de reação do HPA (s) & EKS & 20 & 24.7 & 5.9 & 0.24 & 24.5 & 21.8--27.2 & [22.1, 27.3]\\
\addlinespace
Tempo até o regime permanente (s) & Minikube & 20 & 197.6 & 23.1 & 0.12 & 204.5 & 181.2--218.0 & [187.6, 207.1]\\
Tempo até o regime permanente (s) & EKS & 20 & 83.0 & 46.7 & 0.56 & 69.5 & 57.8--83.2 & [65.8, 105.8]\\
\addlinespace
Réplicas no regime permanente & Minikube & 20 & 6.0 & 0.4 & 0.07 & 6.0 & 6.0--6.0 & [5.9, 6.2]\\
Réplicas no regime permanente & EKS & 20 & 6.7 & 0.6 & 0.09 & 7.0 & 6.0--7.0 & [6.4, 6.9]\\
\addlinespace
CPU de pico observada (\%) & Minikube & 20 & 176.5 & 37.5 & 0.21 & 167.0 & 159.8--179.2 & [161.4, 193.6]\\
CPU de pico observada (\%) & EKS & 20 & 236.2 & 21.5 & 0.09 & 244.0 & 232.0--250.0 & [226.1, 244.2]\\
\addlinespace
Partida do gerador de carga (s) & Minikube & 20 & 2.4 & 1.0 & 0.41 & 2.0 & 2.0--3.0 & [2.0, 2.8]\\
Partida do gerador de carga (s) & EKS & 20 & 1.8 & 0.4 & 0.25 & 2.0 & 1.8--2.0 & [1.6, 1.9]\\
\addlinespace
Retorno ao mínimo após a carga (s) & Minikube & 3 & 453.0 & 6.9 & 0.02 & 449.0 & 449.0--455.0 & [449.0, 461.0]\\
Retorno ao mínimo após a carga (s) & EKS & 3 & 388.7 & 10.7 & 0.03 & 383.0 & 382.5--392.0 & [382.0, 401.0]\\
\addlinespace
Latência de leitura do \emph{API server} (ms) & Minikube & 20 & 108.8 & 3.4 & 0.03 & 108.0 & 106.8--111.2 & [107.4, 110.3]\\
Latência de leitura do \emph{API server} (ms) & EKS & 20 & 547.4 & 53.6 & 0.10 & 519.5 & 516.2--550.5 & [527.0, 571.7]\\
\addlinespace
Latência de escrita no \emph{control plane} (ms) & Minikube & 20 & 464.8 & 575.5 & 1.24 & 197.5 & 147.5--617.8 & [258.4, 743.6]\\
Latência de escrita no \emph{control plane} (ms) & EKS & 20 & 555.2 & 57.1 & 0.10 & 525.5 & 521.5--559.5 & [532.5, 581.0]\\
\addlinespace
Parcela de armazenamento em ms (escrita $-$ leitura) & Minikube & 20 & 355.9 & 576.1 & 1.62 & 92.0 & 33.2--511.5 & [149.7, 630.2]\\
Parcela de armazenamento em ms (escrita $-$ leitura) & EKS & 20 & 7.8 & 14.2 & 1.81 & 6.0 & 2.8--8.0 & [3.0, 14.8]\\
\addlinespace
\bottomrule
\end{tabular}
\end{table*}

% Tabela 2: testes de hipotese.
\begin{table*}[!t]
\centering
\caption{Testes de hipótese, Minikube vs.\ EKS.}
\label{tab:testes}
\footnotesize
\begin{tabular}{@{}llll@{}}
\toprule
\textbf{Teste} & \textbf{Estatística} & $p$ & \textbf{Tamanho de efeito}\\
\midrule
Mann--Whitney $U$ (tempo de reação) & $U=52.0$, $z=-3.99$ & 0.0001 & $\delta=0.740$ (grande)\\
Permutação (diferença de medianas) & $\Delta=9.0$\,s & 0.0004 & --\\
Permutação (razão de dispersões) & razão de MAD $=1.33$ & 0.4098 & --\\
\bottomrule
\end{tabular}
\end{table*}


\subsection{A distribuição do tempo de reação (QP1)}

O tempo de reação no ambiente local teve média de 40,8\,s
($\mathrm{DP}=18{,}5$; mediana 33,5; IIQ 30,8--41,5; faixa 19--87; IC de 95\,\%
da média $[33{,}5;\ 49{,}2]$). No ambiente gerenciado, média de 24,7\,s
($\mathrm{DP}=5{,}9$; mediana 24,5; IIQ 21,8--27,2; faixa 14--36;
IC $[22{,}1;\ 27{,}3]$). As distribuições empíricas aparecem na
Fig.~\ref{fig:ecdf}: a curva do EKS está deslocada à esquerda e é
marcadamente mais íngreme.

A diferença é estatisticamente significativa e o efeito é grande. Um teste de
Mann--Whitney bilateral dá $U=52{,}0$, $z=-3{,}99$, $p=1\times10^{-4}$; um teste
de permutação sobre a diferença de medianas dá diferença de 9,0\,s com
$p=4\times10^{-4}$. O $\delta$ de Cliff é $0{,}740$ e $A_{12}=0{,}870$: uma
repetição sorteada do ambiente local é mais lenta que uma sorteada do ambiente
gerenciado em 87\,\% de todos os pares possíveis.

Vale dizer isso sem rodeios, porque contradiz a conclusão do nosso próprio estudo
anterior. \emph{Existe} um efeito de plataforma sobre o tempo de reação, e o
ambiente gerenciado é o mais rápido. O que faltava ao estudo anterior não era um
argumento melhor, era amostra: com uma observação por ambiente, um efeito desse
tamanho simplesmente não é resolvível contra tanto ruído.

\subsection{Dispersão (QP2)}

O ambiente local é mais variável por toda medida descritiva --- $\mathrm{DP}$ de
18,5\,s contra 5,9\,s, coeficiente de variação de 0,45 contra 0,24, amplitude de
68\,s contra 22\,s. O teste formal, entretanto, não sustenta a afirmação: um
teste de permutação sobre a razão dos desvios absolutos medianos resulta em razão
de 1,33 com $p=0{,}41$.

A discrepância é informativa, e não contraditória. O MAD é uma estatística
robusta, e a distribuição local não é uniformemente espalhada: dezesseis de suas
vinte observações caem entre 19\,s e 43\,s, faixa comparável à do ambiente
gerenciado, enquanto quatro observações (62, 73, 76 e 87\,s) formam uma longa
cauda à direita. A dispersão robusta é, portanto, semelhante; o que difere é a
presença da cauda. Concluímos que a QP2 é respondida negativamente como
formulada --- não conseguimos estabelecer diferença na dispersão \emph{típica}
--- registrando que o ambiente local produziu valores extremos que o gerenciado
nunca produziu.

\subsection{O que as repetições \emph{não} explicam}

Como a campanha registrou covariáveis junto de cada repetição, as fontes dessa
variabilidade podem ser testadas em vez de conjecturadas. Cinco explicações
candidatas foram examinadas por correlação de Spearman contra o tempo de reação,
com $p$-valores por permutação. No ambiente local, \emph{nenhuma} delas é
significativa: latência de escrita no \emph{control plane} ($\rho=-0{,}33$,
$p=0{,}16$), sua parcela de armazenamento ($\rho=-0{,}33$, $p=0{,}15$), partida
do gerador de carga ($\rho=-0{,}33$, $p=0{,}16$), CPU de pico ($\rho=0{,}08$,
$p=0{,}72$) e índice da repetição, que revelaria deriva ao longo da campanha
($\rho=0{,}28$, $p=0{,}22$).

Esse resultado nulo é o achado mais útil da subseção. O \emph{control plane}
local estava demonstravelmente insalubre --- sua latência de escrita teve
coeficiente de variação de 1,24 e variou de 114\,ms a 2534\,ms, contra uma
latência de leitura que se manteve entre 104 e 115\,ms --- e ainda assim essa
instabilidade não prediz quais repetições reagiram devagar. A variabilidade do
tempo de reação, portanto, não é sintoma de um ambiente degradado que uma máquina
melhor removeria.

O mecanismo compatível com essas observações é a \emph{fase}. Como descreve a
Seção~\ref{sec:latencia}, um pico de CPU precisa primeiro ser amostrado pelo
kubelet, depois coletado pelo \emph{Metrics Server} em seu próprio período, e
depois lido pelo HPA em seu período de sincronização. O início da carga cai em um
ponto arbitrário desses ciclos, e esse ponto sozinho pode responder por dezenas
de segundos. A fase não é propriedade da plataforma, não é observável de fora e
não é reprodutível --- o que é precisamente a razão pela qual não se pode confiar
a uma execução única a medição dessa grandeza.

No ambiente gerenciado, duas covariáveis alcançam significância, e ambas
sobrevivem a uma correção de Bonferroni para os dez testes realizados. A CPU de
pico correlaciona positivamente com o tempo de reação ($\rho=0{,}687$,
$p=1{,}4\times10^{-3}$); interpretamos isso como causalidade invertida, e não
como explicação, já que um controlador que decide mais tarde permite à utilização
subir mais antes da primeira ampliação, tornando um pico alto
\emph{consequência} de uma reação lenta. A latência de escrita correlaciona
negativamente ($\rho=-0{,}767$, $p=4\times10^{-4}$), o que não tem leitura causal
que nos convença; o mecanismo plausível é que a sonda de latência roda
imediatamente antes da carga e sua própria duração desloca a fase em que a carga
começa. Que o instrumento perturbe a grandeza que mede é uma limitação que
registramos na Seção~\ref{sec:ameacas}, não um achado.

\subsection{Trajetória, convergência e o resultado lógico (QP4)}

A Fig.~\ref{fig:trajetoria} mostra as trajetórias de réplicas e a
Fig.~\ref{fig:cpu} a utilização sobre a qual o controlador realimentava. Ambos
os ambientes convergem e se estabilizam na meta configurada, reproduzindo o
resultado qualitativo do estudo anterior. Os detalhes quantitativos, porém,
diferem mais do que aquele estudo poderia ter detectado.

O ambiente gerenciado assentou em mais réplicas: mediana 7 (média 6,7; faixa
6--8) contra mediana 6 (média 6,0; faixa 5--7) no local. Isso decorre da
capacidade, e não do controlador: o nó local oferece 4000\,m de CPU alocável em
um único nó, enquanto cada nó de nuvem oferece 1930\,m, e o mesmo gerador de
carga sequencial produziu utilização de pico mais alta na nuvem (236\,\% contra
176\,\% do \emph{request}). A convergência também foi bem mais rápida na nuvem
--- 83,0\,s até o regime permanente contra 197,6\,s no local --- embora com
espalhamento muito maior (CV de 0,56 contra 0,12).

A descida, medida em três repetições por ambiente, levou 453,0\,s (DP 6,9) no
local e 388,7\,s (DP 10,7) na nuvem. Ambas são dominadas pela janela de
estabilização de 300\,s configurada no manifesto, e ambas são notavelmente
consistentes: a descida é governada por política declarada, e política se
reproduz onde temporização não se reproduz.

\subsection{Latência do lado do cliente e sua interpretação}

Os dois valores de latência da Tabela~\ref{tab:descritivas} precisam ser lidos em
par. A latência de leitura do \emph{API server}, que não toca o etcd, foi de
108,8\,ms no local e 547,4\,ms contra o EKS --- esta última é tempo de ida e
volta e negociação TLS até \texttt{us-east-1}, não lentidão do \emph{cluster}. A
parcela de armazenamento, obtida subtraindo a leitura da escrita, conta a
história oposta: 355,9\,ms no local (mediana 92\,ms, máximo 2425\,ms) contra
7,8\,ms na nuvem (mediana 6\,ms). O caminho de armazenamento do \emph{control
plane} gerenciado é cerca de duas ordens de magnitude mais rápido e
incomparavelmente mais estável. Medir apenas a latência de escrita levaria à
conclusão de que os dois \emph{control planes} são semelhantes (464,8\,ms contra
555,2\,ms), o que é exatamente o contrário.

% -----------------------------------------------------------------------------
% Figura: trajetoria do numero de replicas. Mediana entre repeticoes, com faixa
% interquartil. A mediana e preferida a media porque o numero de replicas e
% discreto -- uma media fracionaria nao corresponde a nenhum estado observavel.
\begin{figure}[!t]
\centering
\begin{tikzpicture}
\begin{axis}[estiloBase,
  xlabel={Tempo desde o início da carga (s)},
  ylabel={Réplicas desejadas},
  xmin=-10, xmax=250, ymin=0,
  legend pos=south east,
]
\addplot[draw=none, forget plot, name path=mlo] table[x=t, y=des_q1, col sep=comma] {dados/serie-minikube.csv};
\addplot[draw=none, forget plot, name path=mhi] table[x=t, y=des_q3, col sep=comma] {dados/serie-minikube.csv};
\addplot[corLocal!15, forget plot] fill between[of=mlo and mhi];
\addplot[draw=none, forget plot, name path=elo] table[x=t, y=des_q1, col sep=comma] {dados/serie-eks.csv};
\addplot[draw=none, forget plot, name path=ehi] table[x=t, y=des_q3, col sep=comma] {dados/serie-eks.csv};
\addplot[corNuvem!15, forget plot] fill between[of=elo and ehi];
\addplot[corLocal, const plot] table[x=t, y=des_mediana, col sep=comma] {dados/serie-minikube.csv};
\addlegendentry{Minikube}
\addplot[corNuvem, const plot] table[x=t, y=des_mediana, col sep=comma] {dados/serie-eks.csv};
\addlegendentry{Amazon EKS}
\end{axis}
\end{tikzpicture}
\caption{Trajetória de réplicas ao longo das repetições: mediana (linha) e
amplitude interquartil (faixa) do número de réplicas desejadas, alinhadas pelo
instante em que o contêiner do gerador de carga começou a executar.}
\label{fig:trajetoria}
\end{figure}

% -----------------------------------------------------------------------------
% Figura: ECDF do tempo de reacao. A ECDF e a forma honesta de apresentar uma
% amostra pequena: mostra TODAS as observacoes sem impor um modelo, e a
% separacao (ou sobreposicao) das duas curvas responde a QP1 visualmente.
\begin{figure}[!t]
\centering
\begin{tikzpicture}
\begin{axis}[estiloBase, height=5cm,
  xlabel={Tempo de reação do HPA (s)},
  ylabel={Distribuição acumulada empírica},
  ymin=0, ymax=1.02,
  legend pos=south east,
]
\addplot[corLocal, const plot, mark=*, mark size=1.1pt]
  table[x=reacao_s, y=prob_acumulada, col sep=comma] {dados/ecdf-minikube.csv};
\addlegendentry{Minikube}
\addplot[corNuvem, const plot, mark=*, mark size=1.1pt]
  table[x=reacao_s, y=prob_acumulada, col sep=comma] {dados/ecdf-eks.csv};
\addlegendentry{Amazon EKS}
\end{axis}
\end{tikzpicture}
\caption{Distribuição acumulada empírica do tempo de reação do HPA. Todas as
observações da campanha estão representadas; nenhum modelo distribucional é
imposto.}
\label{fig:ecdf}
\end{figure}

% -----------------------------------------------------------------------------
% Figura: valor de cada repeticao, na ordem em que ocorreu. Serve a dois fins:
% expor a dispersao bruta e permitir inspecionar visualmente se ha TENDENCIA ao
% longo da campanha (aquecimento de cache, deriva termica), que violaria a
% hipotese de repeticoes independentes e identicamente distribuidas.
\begin{figure}[!t]
\centering
\begin{tikzpicture}
\begin{axis}[estiloBase, height=5cm,
  xlabel={Repetição (ordem cronológica)},
  ylabel={Tempo de reação (s)},
  xmin=0, ymin=0,
  legend pos=north east,
]
\addplot[corLocal, only marks, mark=*, mark size=1.5pt]
  table[x=run, y=reacao_s, col sep=comma] {dados/reacao-minikube.csv};
\addlegendentry{Minikube}
\addplot[corNuvem, only marks, mark=square*, mark size=1.5pt]
  table[x=run, y=reacao_s, col sep=comma] {dados/reacao-eks.csv};
\addlegendentry{Amazon EKS}
\end{axis}
\end{tikzpicture}
\caption{Tempo de reação de cada repetição em ordem cronológica. Uma tendência ao
longo da abscissa indicaria que as repetições não são identicamente
distribuídas; nenhuma é aparente.}
\label{fig:porrepeticao}
\end{figure}

% -----------------------------------------------------------------------------
% Figura: utilizacao de CPU observada pelo HPA. Mostra o outro lado do controle:
% a variavel que o controlador realimenta, convergindo para a meta de 50%.
\begin{figure}[!t]
\centering
\begin{tikzpicture}
\begin{axis}[estiloBase, height=5cm,
  xlabel={Tempo desde o início da carga (s)},
  ylabel={Utilização de CPU observada (\%)},
  xmin=-10, xmax=250, ymin=0,
  legend pos=north east,
]
\addplot[corLocal] table[x=t, y=cpu_mediana, col sep=comma] {dados/serie-minikube.csv};
\addlegendentry{Minikube}
\addplot[corNuvem] table[x=t, y=cpu_mediana, col sep=comma] {dados/serie-eks.csv};
\addlegendentry{Amazon EKS}
\addplot[corNeutra, dashed, line width=0.7pt, forget plot]
  coordinates {(-10,50) (250,50)};
\node[font=\scriptsize, corNeutra, anchor=south west] at (axis cs:150,52) {meta = 50\,\%};
\end{axis}
\end{tikzpicture}
\caption{Utilização mediana de CPU observada pelo HPA ao longo das repetições. A
variável controlada converge para a meta configurada nos dois ambientes.}
\label{fig:cpu}
\end{figure}

% =============================================================================
\subsection{Quão enganosa é uma execução única?}
\label{sec:execucao-unica}

As distribuições medidas permitem avaliar diretamente o desenho do estudo
anterior. Sorteando uma repetição de cada ambiente --- que é exatamente o que um
experimento de execução única faz --- a probabilidade de cada conclusão possível
é: \textbf{87,0\,\%} de o ambiente gerenciado parecer mais rápido,
\textbf{12,2\,\%} de o local parecer mais rápido e 1,5\,\% de empate exato. Um
experimento de execução única neste sistema, portanto, erra a direção do efeito
em aproximadamente uma vez a cada oito.

Esse número já explica a irreprodutibilidade do estudo anterior, mas as
observações individuais podem ser localizadas com precisão ainda maior. A
Tabela~\ref{tab:percentis} situa cada valor previamente relatado em seu percentil
empírico dentro da distribuição aqui medida.

\begin{table}[!t]
\centering
\caption{Onde as observações únicas do estudo anterior caem dentro das
distribuições medidas nesta campanha.}
\label{tab:percentis}
\footnotesize
\begin{tabular}{@{}llrr@{}}
\toprule
\textbf{Execução anterior} & \textbf{Ambiente} & \textbf{Relatado} & \textbf{Percentil}\\
\midrule
E1 (Andrade)    & Minikube & 62\,s & 82\,\%\\
E1 (Andrade)    & EKS      & 15\,s & 10\,\%\\
E2 (Miqueletti) & Minikube & 33\,s & 45\,\%\\
E2 (Miqueletti) & EKS      & 60\,s & 100\,\%\\
\bottomrule
\end{tabular}
\end{table}

O padrão está completo. A execução E1 por acaso sorteou uma observação local
lenta (percentil 82) junto de uma observação de nuvem rápida (percentil 10), e
assim relatou uma vantagem de quatro vezes para a nuvem --- exagerando um efeito
real, porém muito menor. A execução E2 sorteou uma observação local absolutamente
ordinária (percentil 45) junto de uma observação de nuvem acima de qualquer coisa
que registramos em vinte repetições, e assim inverteu a conclusão. Nenhuma das
duas execuções foi defeituosa; ambas foram sorteios únicos de distribuições que
se sobrepõem.

Uma ressalva se aplica à última linha. O estudo anterior media $T_0$ na submissão
do gerador de carga pelo operador, embutindo o escalonamento e a partida do Pod,
que aqui medimos em 1,8\,s na nuvem. Isso infla ligeiramente os valores antigos e
não explica, por si só, uma observação de 60\,s exceder nosso máximo de 36\,s; a
explicação mais provável é que a instrumentação mais grosseira e manual de E2
tenha superestimado o valor. Trata-se da mesma ameaça de instrumentação
heterogênea que motivou o desenho atual.

% =============================================================================
\section{Discussão}
\label{sec:discussao}

\subsection{O que replica e o que não replica}

Os resultados se separam nitidamente em dois grupos, e a separação não é a que o
estudo anterior propunha.

Replica tudo o que é determinado por declaração. Ambos os ambientes convergiram,
ambos estabilizaram a utilização média na meta configurada de 50\,\%, e ambos
desceram conforme a cláusula \texttt{behavior}, com tempos de descida cujos
desvios-padrão foram de 6,9\,s e 10,7\,s sobre médias de 453\,s e 389\,s ---
variação de cerca de 2\,\%. A Equação~\eqref{eq:hpa} não contém termo de
infraestrutura, e as medições se comportam de acordo. Quem quer saber se um
manifesto \emph{está correto} pode responder isso localmente.

Não replica a temporização, mas a formulação anterior desse fato estava errada. O
tempo de reação não é irreprodutível: é uma variável aleatória com distribuição
estável, e essas distribuições diferem entre ambientes por uma margem grande e
significativa. O que é irreprodutível é um \emph{sorteio único} dela. A distinção
importa porque muda o remédio: a resposta não é um experimento único melhor
controlado, é repetição.

\subsection{Por que o ambiente gerenciado reage mais rápido}

Nossos dados sustentam o efeito, mas não isolam seu mecanismo, e tomamos cuidado
para não afirmar mais do que se pode. O candidato óbvio --- que o caminho de
armazenamento do \emph{control plane} local é duas ordens de magnitude mais
lento --- é contrariado pela análise de correlação: dentro do ambiente local, as
repetições com pior latência de escrita não foram as mais lentas. O que produz o
deslocamento parece agir de forma uniforme, e não através dos episódios de
degradação de armazenamento que capturamos.

Uma explicação plausível, consistente com tudo o que foi medido, é que o
\emph{control plane} gerenciado executa sua cadeia de métricas e sua malha de
controle sobre infraestrutura dedicada e consistentemente provisionada, de modo
que os \emph{períodos} dessas malhas são honrados com precisão. O \emph{control
plane} local compartilha CPU e um caminho de E/S saturado com o hospedeiro, de
modo que uma coleta ou uma sincronização pode escapar de seu período nominal. Um
escorregamento desse tipo deslocaria a distribuição inteira e alongaria sua
cauda, que é o que observamos: quatro repetições locais além de 60\,s contra um
máximo de 36\,s na nuvem. Testar isso adequadamente exigiria instrumentar os
tempos internos da malha do controlador, que o serviço gerenciado não expõe.

\subsection{Implicações práticas}

Três recomendações decorrem diretamente das medições.

Primeira: \textbf{valide a configuração de autoescalamento localmente e meça sua
temporização no ambiente de destino}. A documentação da própria AWS para o
EKS~\cite{awshpa} descreve a configuração, mas nada diz sobre a latência
esperada, que, a julgar por nossos resultados, não poderia de todo modo ser
declarada como um número único. O \emph{cluster} local reproduziu o resultado
lógico com fidelidade e baixo custo. Não reproduziu a latência, e nenhum cuidado
adicional numa execução local única teria revelado isso.

Segunda: \textbf{não calibre parâmetros de temporização a partir de uma
observação única}. Um operador que medisse 62\,s no local e 15\,s na nuvem ---
ambos valores que observamos --- concluiria que a plataforma responde por uma
diferença de quatro vezes. A diferença real de medianas é de 9\,s. Dimensionar
uma janela de estabilização ou um limiar de alerta a partir de uma execução única
corre o risco de codificar um artefato de fase.

Terceira: \textbf{relate distribuições, não valores pontuais}. O custo da
campanha de vinte repetições aqui relatada foi de cerca de três horas de execução
não supervisionada por ambiente e, a preços de tabela, menos de um dólar
americano de consumo de nuvem. Contra isso, o custo de publicar uma conclusão
invertida é considerável, e nosso estudo anterior chegou a uma execução de
distância de fazer exatamente isso.

\subsection{Custo e assimetria operacional}

A campanha de nuvem manteve um \emph{cluster} vivo por cerca de quatro horas e
meia, incluindo um reinício. Reportamos o custo com uma distinção que costuma ser
omitida: o custo \emph{faturado}, consultado após o fato na API de gestão de
custos da conta, foi zero --- a conta usada neste trabalho não é cobrada por uso
de EKS, EC2 ou VPC --- enquanto o custo \emph{a preço de tabela} do mesmo
consumo, às tarifas públicas sob demanda de um \emph{control plane}, duas
instâncias \texttt{t3.small} e um \emph{NAT gateway}, é de aproximadamente
US\$\,0,90. Apenas o segundo é informativo para quem planeja uma campanha
semelhante, e é uma estimativa, não uma medição. A campanha local não custou nada
sob nenhuma das duas contabilidades e levou aproximadamente o mesmo tempo de
parede.

A assimetria que importa não é o preço, mas a forma do modo de falha: o ambiente
local falhou por degradação --- durante a preparação seu \emph{control plane}
tornou-se inacessível, com o etcd reportando operações de \texttt{fdatasync} de
1,9\,s --- e exigiu um guarda de saúde entre repetições para sobreviver a uma
campanha não supervisionada. O ambiente gerenciado nunca precisou disso. Nenhuma
das campanhas perdeu repetições, mas apenas uma delas precisou da rede de
segurança.

\subsection{Generalidade do resultado metodológico}

Nada no argumento da Seção~\ref{sec:execucao-unica} é específico do HPA. Ele
exige apenas que a grandeza de interesse seja uma variável aleatória com
distribuições sobrepostas entre condições, e que o experimento a amostre uma vez
por condição. Isso descreve uma fração grande das comparações de infraestrutura
de nuvem feitas por profissionais, onde replicar é barato, a variabilidade é alta
e comparações de execução única são a norma. A lição transferível é a aritmética:
quando $A_{12}=0{,}87$, uma execução por condição inverte a ordenação uma vez a
cada oito, e o experimento não oferece nenhum sinal de que isso aconteceu.

% =============================================================================
\section{Ameaças à validade}
\label{sec:ameacas}

\subsection{Validade de construção}

\paragraph{O tempo de reação mede a decisão, não o efeito} Nossa métrica primária
termina na decisão de escala do controlador. O usuário só sente alívio quando
novos Pods estão servindo tráfego, o que acrescenta atrasos de escalonamento, de
imagem e de prontidão. Reportamos a latência de partida do gerador e o tempo até
o regime permanente separadamente precisamente porque esses custos de atuação são
onde os ambientes podem legitimamente diferir, mas um leitor interessado em
elasticidade de ponta a ponta deve combinar as duas coisas em vez de ler apenas o
tempo de reação.

\paragraph{Resolução de um segundo nos carimbos do cluster} Tanto
\texttt{startedAt} quanto \texttt{lastScaleTime} são carimbos RFC\,3339 com
granularidade de segundo, de modo que todo tempo de reação carrega até $\pm1$\,s
de quantização, além da quantização muito mais grosseira imposta pelo período de
15\,s do controlador. Isso afeta de forma desprezível as estimativas de posição,
mas infla o número de empates exatos, razão pela qual a correção de empates no
teste de Mann--Whitney não é opcional aqui.

\subsection{Validade interna}

\paragraph{Repetições dentro de uma única sessão} Cada campanha roda como uma
única sessão ininterrupta em um único \emph{cluster}. Efeitos que derivam ao
longo de horas --- aquecimento do cache de imagens, comportamento térmico do
hospedeiro, atividade de vizinhos do lado da AWS --- ficam, portanto, confundidos
com o índice da repetição em vez de serem promediados. Mitigamos isso verificando
um estado de base explícito antes de cada repetição e publicando o tempo de
reação contra a ordem cronológica (Fig.~\ref{fig:porrepeticao}), de modo que uma
tendência ficaria visível em vez de escondida dentro de um agregado.

\paragraph{Descida forçada na maioria das repetições} Apenas as três primeiras
repetições de cada campanha observam a descida natural; as demais forçam o
retorno a uma réplica. As estatísticas de descida, portanto, repousam sobre
$n=3$ por ambiente e são reportadas como secundárias e indicativas, não como
distribuição caracterizada. O próprio ato de forçar não pode contaminar o tempo
de reação da repetição seguinte, porque a verificação do estado de base exige que
a CPU medida tenha decaído antes de a próxima repetição começar.

\paragraph{O ambiente local compartilha seu hospedeiro} O \emph{cluster} Minikube
disputa CPU com o sistema operacional do hospedeiro e com o próprio cliente de
medição. Isso não é um defeito a ser eliminado por engenharia: é a condição real
de um \emph{cluster} local de desenvolvimento, e faz parte daquilo de que a
comparação trata. Significa, porém, que os resultados locais caracterizam
\emph{este} hospedeiro, e não Kubernetes local em geral.

\subsection{Validade externa}

\paragraph{Uma carga, um gerador fraco} Um único laço sequencial de
\texttt{busybox} nunca levou nenhum dos \emph{clusters} ao teto de dez réplicas.
A campanha caracteriza o comportamento do controlador em regime de convergência,
não os limites de capacidade de nenhum dos ambientes, e nada diz sobre como o
tempo de reação se comporta quando o \emph{cluster} está próximo da saturação ou
quando Pods ficam \texttt{Pending} por falta de capacidade de nó.

\paragraph{Uma topologia, uma região, um tipo de instância} O braço de nuvem são
duas instâncias \texttt{t3.small} em \texttt{us-east-1}. Instâncias burstáveis
têm dinâmica de créditos de CPU que famílias maiores não têm, e uma única região,
em uma única conta, em um único dia, é uma fatia estreita de um serviço
gerenciado. Escolhemos essa topologia porque é a que o estudo anterior usou, e a
comparabilidade com E1 e E2 era o ponto.

\paragraph{Ajuste padrão do controlador} O período de sincronização do HPA e a
resolução do \emph{Metrics Server} foram mantidos em seus valores padrão, e são
os termos dominantes no orçamento de latência da Seção~\ref{sec:latencia}.
Conclusões sobre tempo de reação são conclusões sobre um HPA \emph{com ajuste
padrão}; um operador que reduza o período de sincronização estaria medindo um
sistema diferente.

\subsection{Validade de conclusão}

\paragraph{Tamanho da amostra} Vinte repetições por ambiente sustentam uma
descrição distribucional e um teste não paramétrico, mas estimam mal as caudas.
As afirmações deste artigo são sobre tendência central, dispersão e ordenação
estocástica; não fazemos nenhuma afirmação sobre percentis extremos.

\paragraph{Comparações múltiplas} Várias métricas são comparadas entre ambientes.
A análise confirmatória é a declarada em QP1 e QP2 --- posição e dispersão do
tempo de reação; toda outra comparação é exploratória e deve ser lida como tal,
e não como achado significativo independente.

% =============================================================================
\section{Conclusão}
\label{sec:conclusao}

Este artigo reexaminou uma questão que nosso próprio trabalho anterior havia
respondido incorretamente: se o tempo de reação do \emph{Horizontal Pod
Autoscaler} do Kubernetes é uma propriedade da plataforma. Aquele trabalho
implantou manifestos idênticos em Minikube e Amazon EKS, executou cada
implantação uma vez por autor, e constatou que a ordenação dos dois ambientes se
invertia entre execuções --- de onde concluiu que a diferença não era atribuível
à plataforma. Aqui mantivemos a implantação, a carga e o protocolo inalterados,
mudamos o desenho de uma para vinte execuções por ambiente, e derivamos a métrica
primária de carimbos de tempo atribuídos pelo \emph{cluster} em vez de pelo
observador.

Respondendo às questões de pesquisa: \textbf{(QP1)} o tempo de reação é uma
variável aleatória cujas distribuições diferem significativamente entre ambientes
--- média de 40,8\,s no local contra 24,7\,s na nuvem, $U=52{,}0$,
$p=1\times10^{-4}$, $\delta$ de Cliff $=0{,}74$ --- de modo que existe um efeito
de plataforma, e ele favorece o serviço gerenciado; \textbf{(QP2)} o ambiente
local produziu uma longa cauda à direita que o gerenciado nunca produziu, mas um
teste robusto de dispersão não estabelece diferença no espalhamento típico;
\textbf{(QP3)} uma execução única por ambiente reproduz a ordenação correta em
apenas 87\,\% das vezes, e as duas observações anteriores situam-se nos percentis
82 e 10 em uma execução e 45 e 100 na outra, o que explica integralmente a
inversão; \textbf{(QP4)} todo resultado fixado por declaração --- convergência,
meta de utilização, descida assimétrica --- replicou entre ambientes com variação
de poucos por cento.

A correção à nossa conclusão anterior merece ser enunciada com precisão, porque a
formulação anterior não era apenas imprecisa, era enganosa. Não é o caso de a
latência de autoescalamento não se relacionar com a plataforma. É o caso de a
diferença ser menor do que um único par de observações sugere, e de o ruído em
torno dela ser grande o bastante para inverter a ordenação com frequência. Um
resultado negativo obtido com $n=1$ por condição carrega essencialmente nenhuma
informação, e apresentá-lo como achado --- como fizemos --- afirma mais do que os
dados sustentavam.

Os resultados de instrumentação são transferíveis além deste experimento. Derivar
a latência de \texttt{lastScaleTime} e do \texttt{startedAt} do contêiner remove
inteiramente o relógio do observador da medição; medir as latências de leitura e
de escrita do \emph{API server} como um par isola o caminho de armazenamento do
\emph{control plane} da distância de rede, que a latência de escrita bruta
confunde tão gravemente que chega a inverter a ordenação; e substituir o
\emph{exec credential plugin} do EKS por um \emph{token} renovado foi necessário
para amostrar os dois ambientes na mesma taxa.

Trabalhos futuros seguem as limitações, e não os resultados. A extensão mais
valiosa instrumentaria os tempos internos da malha do controlador para testar a
hipótese de fase diretamente, em vez de inferi-la da ausência de correlação com
todas as covariáveis que conseguimos observar de fora. Um gerador de carga
paralelo caracterizaria o comportamento próximo ao teto de réplicas, onde o
\emph{Cluster Autoscaler} se torna necessário e o tempo de reação pode se
comportar de outra forma. Uma campanha que se estendesse por vários dias e
famílias de instância separaria efeitos da plataforma gerenciada de efeitos de
uma conta em uma tarde. Por fim, o mesmo desenho de medidas repetidas aplicado ao
KEDA e ao autoescalamento por métricas customizadas estabeleceria se a razão
entre reprodutibilidade declarativa e variabilidade temporal aqui encontrada é
propriedade geral das malhas de controle do Kubernetes.

\section*{Reprodutibilidade}

Os manifestos, os \emph{scripts} de campanha, as séries temporais brutas de todas
as quarenta repetições, os registros de eventos do controlador, a análise
estatística e o código-fonte deste artigo estão disponíveis no repositório do
projeto. A análise usa apenas a biblioteca padrão do Python e uma semente
aleatória fixa, de modo que cada figura e cada número das tabelas se regeneram a
partir dos dados brutos com um único comando. As figuras deste artigo são
produzidas pelo \textsc{pgfplots} lendo diretamente a saída da análise, sem
arquivos de imagem intermediários.

% =============================================================================
\begin{thebibliography}{9}

\bibitem{nist}
P. Mell and T. Grance, \emph{The NIST Definition of Cloud Computing}, NIST
Special Publication 800-145. Gaithersburg, MD: National Institute of Standards
and Technology, 2011.

\bibitem{k8shpa}
Kubernetes, \emph{HorizontalPodAutoscaler Walkthrough}. [Online]. Available:
\url{https://kubernetes.io/docs/tasks/run-application/horizontal-pod-autoscale-walkthrough/}

\bibitem{k8salgo}
Kubernetes, \emph{Horizontal Pod Autoscaling --- Algorithm Details}. [Online].
Available:
\url{https://kubernetes.io/docs/tasks/run-application/horizontal-pod-autoscale/}

\bibitem{metricsserver}
Kubernetes SIG, \emph{Metrics Server}. [Online]. Available:
\url{https://github.com/kubernetes-sigs/metrics-server}

\bibitem{awshpa}
Amazon Web Services, \emph{Scale Pod Deployments with Horizontal Pod
Autoscaler}, Amazon EKS User Guide. [Online]. Available:
\url{https://docs.aws.amazon.com/eks/latest/userguide/horizontal-pod-autoscaler.html}

\bibitem{eksctl}
\emph{eksctl --- The Official CLI for Amazon EKS}. [Online]. Available:
\url{https://eksctl.io/}

\bibitem{mytkowicz}
T. Mytkowicz, A. Diwan, M. Hauswirth, and P. F. Sweeney, ``Producing wrong data
without doing anything obviously wrong!'' in \emph{Proc. ASPLOS}, 2009,
pp. 265--276.

\bibitem{hoefler}
T. Hoefler and R. Belli, ``Scientific benchmarking of parallel computing
systems,'' in \emph{Proc. SC}, 2015, pp. 1--12.

\bibitem{arcuri}
A. Arcuri and L. Briand, ``A Hitchhiker's guide to statistical tests for
assessing randomized algorithms in software engineering,'' \emph{Softw. Test.
Verif. Reliab.}, vol. 24, no. 3, pp. 219--250, 2014.

\end{thebibliography}

\end{document}
